% TOP_PROBLEM: {"id": 30007051, "title": "The ordinary-double-point normalized-volume gap", "category_id": 5, "category": {"id": 5, "name": "algebraic_geometry", "display_name": "Algebraic Geometry", "description": "Geometric objects defined by polynomial equations.", "slug": "algebraic-geometry", "order_index": 5, "created_at": "2026-07-31T15:26:25.671Z"}, "status": "open"}
% FIRST_POSTED: 2026-10-06
% ATTEMPT_STATUS: solved
% OPENAI_CLAIM: {"schema_version":1,"source_repo":"https://github.com/openai/math","source_commit":"adc7f1241b42e322a6451854ab7e4b4c146bf78a","release_date":"2026-10-06","match":"full","resolution":"proved","independently_reviewed":false,"second_pass":"agreed","families":[{"family":"037","title":"The ordinary-double-point volume gap","subject":"Algebraic and complex geometry","manuscripts":[{"title":"The ordinary-double-point gap in every dimension","pdf_path":"preprints/The-ordinary-double-point-gap-in-every-dimension-September-24-2026/paper.pdf","readme_path":"preprints/The-ordinary-double-point-gap-in-every-dimension-September-24-2026/README.md","date":"2026-09-24","theorem_ref":"thm:main (sections/01-introduction.tex:33)","statement_sources":[{"path":"preprints/The-ordinary-double-point-gap-in-every-dimension-September-24-2026/build/sections/01-introduction.tex","line":33,"label":"thm:main"}]}],"lean":null,"reasoning_trace":null}]}
\documentclass[11pt]{article}
\usepackage[margin=25mm]{geometry}
\usepackage{fontspec}
\setmainfont{Latin Modern Roman}
\usepackage{amsmath,amssymb,mathtools}
\usepackage{unicode-math}
\setmathfont{Latin Modern Math}
\usepackage{enumitem}
\usepackage{xurl,hyperref,bookmark}
\hypersetup{colorlinks=true,urlcolor=blue}
\setcounter{secnumdepth}{0}
\setlength{\parindent}{0pt}
\setlength{\parskip}{5pt}
\setlength{\emergencystretch}{3em}
\begin{document}
\section*{The ordinary-double-point normalized-volume gap}

\subsection{OpenAI's claim}
Every singular boundary-zero complex klt germ of dimension $n$ is claimed to have normalized volume at most $2(n-1)^n$, with equality exactly for an analytic ordinary double point. The cited manuscript argues for that result; its relation to this problem is specified below.

\subsection{Mathematical statement of the claimed result}
\paragraph{Mathematical claim 1.} Let $x$ be a closed point that is singular on a normal algebraic complex $n$-fold, $n\ge2$, with rational Cartier $K_X$ and boundary-zero klt singularity near $x$. Define normalized volume by
\[
\widehat{\operatorname{vol}}(x,X)=\inf_{v\text{ centered at }x}A_X(v)^n\operatorname{vol}(v),
\]
where $A_X$ is log discrepancy and $\operatorname{vol}(v)=\limsup_{t\to\infty}n!\,\operatorname{length}(\mathcal O_{X,x}/\{f:v(f)\ge t\})/t^n$. OpenAI claims $\widehat{\operatorname{vol}}(x,X)\le2(n-1)^n$, with equality exactly when the analytic germ is the ordinary double point $\{z_0^2+\cdots+z_n^2=0\}\subset\mathbb C^{n+1}$. All singular boundary-zero klt germs are quantified, not only hypersurfaces.

\subsection{Outline of the argument}
Manuscript 1 (family 037, 2026-09-24): The four-dimensional companion result anchors an induction using minimizing cones and integral approximations of their Reeb gradings. Cyclic stabilizer characters are retained through orbifold quotients, and same-dimensional singular slices compare the remaining volume bounds and equality cases.

\subsection{Scope relative to this problem}
Manuscript 1 (family 037, 2026-09-24): The problem selects the mathematical statement summarized here, including its stated hypotheses and parameter range. The cited theorem gives a claimed proof of that assertion in full.

\subsection{Formal verification}
Family 037 has no Lean scope page in this snapshot; no formalization of these manuscript results is listed.

This site has not executed the Lean files or checked their dependencies. These descriptions report the scope of the linked declarations, rather than a new verification of the proofs.

\subsection{Status caveat}
OpenAI attributes these results to an internal model. Its release notice cautions that results lacking full formalization can contain mistakes. This site has not independently assessed the proofs; a solved label records the scope of the reported claim.

\subsection{Sources}
\begin{itemize}
\item OpenAI, \emph{The ordinary-double-point gap in every dimension}, 2026-09-24; citation key \texttt{OAI:The-ordinary-double-point-gap-in-every-dimension-September-24-2026}. \href{https://github.com/openai/math/blob/adc7f1241b42e322a6451854ab7e4b4c146bf78a/preprints/The-ordinary-double-point-gap-in-every-dimension-September-24-2026/paper.pdf}{PDF (GitHub)}; \href{https://github.com/openai/math/raw/adc7f1241b42e322a6451854ab7e4b4c146bf78a/preprints/The-ordinary-double-point-gap-in-every-dimension-September-24-2026/paper.pdf}{PDF (direct)}; \href{https://github.com/openai/math/blob/adc7f1241b42e322a6451854ab7e4b4c146bf78a/preprints/The-ordinary-double-point-gap-in-every-dimension-September-24-2026/README.md}{Manuscript page and citation}; \href{https://github.com/openai/math/blob/main/preprints/The-ordinary-double-point-gap-in-every-dimension-September-24-2026/paper.pdf}{Latest version}; Statement source: \href{https://github.com/openai/math/blob/adc7f1241b42e322a6451854ab7e4b4c146bf78a/preprints/The-ordinary-double-point-gap-in-every-dimension-September-24-2026/build/sections/01-introduction.tex\#L33}{thm:main}.
\item \href{https://github.com/openai/math/blob/adc7f1241b42e322a6451854ab7e4b4c146bf78a/CONTENTS.md}{OpenAI catalogue entry}.
\end{itemize}
\end{document}
